\section{Discussion}
\label{sec:discussion}

The construction answers a compatibility question. It shows that a single stationary, isotropic, simple process can have exact reciprocal full correlations at orders two, three and four on open finite domains. The correlation identities include every allocation of tuple points to different parents. This is the main point of the proof: lower-order residuals must be chosen so that the full densities, rather than only the same-parent contributions, have the required form.

The domain restrictions have separate roles. The finite maximum distance gives finite total integrals for the reciprocal kernels. The positive minimum altitude for tetrahedra keeps the reciprocal-volume kernel integrable in the omitted-vertex direction. The reciprocal-distance kernel is already locally integrable at zero in three dimensions; the lower pair-distance cutoff specifies the resolved support and, together with the triangle altitude cutoff, gives support nesting. These are sufficient features of the construction. The proof does not claim that each cutoff is necessary for every point-process architecture.

The explicit witness has extremely low intensity. This is not an incidental numerical defect in the proof: small amplitudes make the subtractive kernel corrections uniformly small, which certifies positivity. We have not optimised intensity or shown that another finite-cluster law cannot do better. The model also has no time dynamics, magnitude marks, fitted fault geometry or observation process. The finite-range independence is a useful structural property of this witness, but it does not supply these missing earthquake features.

The numerical validation in Section~\ref{sec:simulation} checks a particular sampler and a family of finite-bin estimators; it does not establish the continuum identities, which were proved in Theorem~\ref{thm:main}. The $K$-function and pooled pair and triangle estimates reproduce their targets at useful precision. Fourth order does not: 126 admissible tetrahedra leave every pooled volume bin underpowered. This is a computational consequence of the selected sparse witness, and it also illustrates why an exact existence theorem need not produce an efficient numerical or empirical model. Conversely, exact continuum identities do not determine the finite-catalogue ratios used in the historical studies. Section~\ref{sec:earthquake-comparison} records the observation model still needed for that comparison.

Several questions remain. Can one raise the intensity while preserving all three exact laws and finite local moments? Which broader shape domains admit the same construction? What additional constraints on count distributions or temporal dependence would make a stationary spatial model useful for earthquake catalogues? The present proof supplies one exact starting point for these questions, not their answers.
